\subsection{Adapter: Windowed Protection Control}
\label{sec:kv-runtime}
\label{sec:capacity-allocation}

The Adapter (4) adjusts protection to failure risk and foreground
memory demand. At each control-window boundary, it publishes a checkpoint interval $Z$ in
generated tokens, a target Peer depth $L$, and a total Peer budget $B_p$
for each target DP. The Controller coordinates Adapter reports so all
requests of the same deployed model use the same $Z$ and $L$, including
new arrivals. These settings remain fixed within the window; physical
allocation still checks current free space. The KV Master makes request
and block decisions under these settings, while Connectors move bytes.

\paragraph{Adjust checkpoint frequency to failure risk.}
Failure risk comes from a smoothed history of failures per unit of
device operating time in the relevant failure domain. A configured
prior is used when observations are insufficient; a fault-free control
window does not imply zero risk. Higher risk favors shorter checkpoint
intervals, balancing normal-operation overhead against lost work, as in
classical checkpointing~\cite{daly2006checkpoint}. Here, $Z$ advances by
one level in a configured, block-aligned list. The Adapter shortens $Z$
above an upper risk threshold only when copy backlog, commit lag, and
foreground interference stay within configured guards. Below a lower
threshold it lengthens $Z$; otherwise it retains the current value.
At either end of the allowed list, it stays at that bound. Thresholds,
copy guards, candidate depths, and switching margins are deployment settings.
Incremental asynchronous copying requires its own calibration, rather
than directly applying a blocking-checkpoint formula.

$Z$ triggers a Master decision, not a completed replica. The Master
offers sealed state after prefill and revisits an active request after
each additional $Z$ generated tokens. Publishing a new interval does
not reset that request's progress counter. The actual replay tail is
the known token prefix beyond the last valid committed checkpoint;
copy queuing can make it longer than $Z$. Tier R need not checkpoint.
The Adapter therefore records commit lag as well as decision frequency.

\paragraph{Reserve memory for foreground inference.}
For each DP, the next budget is the managed KV pool minus estimated
foreground KV demand and admission headroom, clamped between zero and
half of that pool. The pool excludes weights and fixed runtime state,
but includes existing Peer copies. Thus $B_p$ is a total target, not
newly available space. Rising foreground memory demand reduces this
target; sustained relief permits expansion. KV watermarks and
memory-related admission stalls inform demand, rather than GPU compute
utilization alone.

\paragraph{Choose one depth for the next window.}
A shorter interval usually leaves less tail computation to overlap
DRAM loading, making additional Peer layers useful. A longer interval
can reach diminishing returns at a smaller depth. These are tendencies:
failures can occur anywhere between checkpoints, and larger $L$ may
leave fewer requests protected under the same budget.

The Adapter compares a small configured set of depths, including zero
and the current $L$ (Algorithm~\ref{alg:adapter}). Its recovery table
covers checkpoint intervals, context lengths, copy bandwidth, replay
costs, and concurrency. Entries average over time-sampled checkpoint
phases from backup/replay profiles at each interval, including commit
lag. This assumes failure risk is independent of checkpoint phase,
rather than a failure at exactly $Z$ tokens. Recent workload statistics
select and weight the entries.
For a candidate depth, a preview of the Master's admission policy
accounts for shared objects, referenced copies, pending transfers, and
the next-window copy limit; it credits only expected committed coverage.
The score adds predicted foreground delay from memory displacement and
backup work to recovery delay not absorbed by request slack, weighted
by each source's failure probability in the next window. This sum of request delays is a
selection heuristic, not an additive estimate of workflow makespan.
The lowest score wins, with smaller depths breaking ties. A change
requires an improvement beyond the configured uncertainty margin,
compared with keeping the current depth under the new interval and budgets;
missing profiles retain the current depth. This avoids deriving a
separate $L$ for every request or assuming that higher risk always
permits a larger depth.

\begin{figure}[t]
  \raggedright
  \refstepcounter{designalgorithm}\label{alg:adapter}
  \noindent\textbf{Algorithm \thedesignalgorithm: Windowed protection update}
  \par\smallskip
  \noindent Input: current settings, risk and load reports, recovery table.
  Output: next-window $Z$, common $L$, and budgets $B_p$.
  \par\smallskip
  \begin{tabularx}{\columnwidth}{@{}r@{\quad}X@{}}
    \toprule
    1 & Snapshot live requests, references, and copy progress \\
    2 & Estimate failure risk from the longer-term history \\
    3 & \textbf{for each} target DP $p$: \\
    4 & \quad$B_p\gets$ KV pool $-$ foreground demand $-$ headroom \\
    5 & \quad Clamp $B_p$ to $[0,\ \text{half of KV pool}]$ \\
    6 & $Z_{\mathrm{next}}\gets Z$; $L_{\mathrm{next}}\gets L$ \\
    7 & \textbf{if} risk is high and copy/foreground guards allow: \\
    8 & \quad$Z_{\mathrm{next}}\gets$ next smaller allowed interval \\
    9 & \textbf{else if} risk is low: \\
    10 & \quad$Z_{\mathrm{next}}\gets$ next larger allowed interval \\
    11 & \textbf{if} profiles cover $Z_{\mathrm{next}}$ and current load: \\
    12 & \quad\textbf{for each} candidate depth, including $0$ and $L$: \\
    13 & \qquad Preview admissions and committed coverage \\
    14 & \qquad Score $\gets$ estimated foreground delay \\
    15 & \qquad Add risk-weighted recovery delay beyond slack \\
    16 & \quad Pick minimum score; break ties by smaller depth \\
    17 & \quad \textbf{if} materially better: $L_{\mathrm{next}}\gets$ chosen depth \\
    18 & Publish $Z_{\mathrm{next}},L_{\mathrm{next}},\{B_p\}$ for the next window \\
    19 & Pause excess admission; reclaim only safe copies \\
    \bottomrule
  \end{tabularx}
  \Description{The Adapter reserves foreground memory, adjusts the
    checkpoint interval using failure risk and copy guards, evaluates
    common layer depths, and publishes one set of window settings.
    Reclamation preserves active references and committed recovery state.}
\end{figure}

% Reference-aware execution of the window targets, part of section 3.6.
\paragraph{Admit and reclaim without evicting live references.}
The Master offers hot blocks from the first $L$ layers in slack order,
using stable request IDs to break ties. Existing copies can be shared;
new copies require both budget and physical space, reserved before
transfer. Under pressure, only zero-reference objects with no in-flight
reader or writer and valid DRAM backing enter the LRU eviction queue.
If free space and safe victims are insufficient, the Peer increment is
denied and the request retains its DRAM path. Normal inference does not
wait for a Peer grant, and priority does not revoke a live reference.

A successful grant holds a request reference until that LLM request
ends, including a response that hands control to an external tool.
Completion decrements shared references; zero references make a copy
eligible for reclamation, not immediate deletion. Growing or changed
hot sets require fresh admission. Previously granted objects remain
protected, so the current candidate count is a target rather than a
guarantee of complete Peer coverage or a bound on retained old copies.

When the budget shrinks, new allocations stop above the target and
safe zero-reference copies are reclaimed first. Referenced copies drain
as requests finish; reducing $L$ likewise does not revoke their pins.
Physical usage can therefore exceed a newly reduced target, and memory
is returned to foreground inference only after release completes.
Increasing $L$ queues additional candidates subject to the same checks.
Copying, committed, and reclaiming bytes remain charged until physical
release. Window targets never imply instant copies or instant free space.


\paragraph{Commit and recover by actual readiness.}
Peer and fault-independent DRAM paths commit separately, after transfer
completion and model/layout/shard validation. Low-priority copying is
rate-limited (default 20\% of link capacity), with target HBM bandwidth
also accounted for. Recovery reads committed objects, preloads hot
blocks in layer order, and gives newly revealed misses priority over
speculative copies. Selector state arrives before selection; each layer
waits for all required inputs. Actual DMA granularity can limit copy
preemption. Tier R rebuilds history. Placement changes neither selector
outputs nor known replay tokens and adds no attention approximation;
compatible weights, layout, and operators reconstruct the same logical
KV state, without promising identical future stochastic samples across
floating-point schedules.

\begin{figure}[t]
  \centering
  \includegraphics[width=\columnwidth]{figures/design-single/design-layer-pipeline.pdf}
  \caption{Modeled recovery pipelines: MiniCPM4-8B, 32 layers, 32K
    history, 64-token known tail; $C=2.87$\,ms/layer (measured) and
    $D=5.37$\,ms at an assumed 6.25\,GB/s. In (d), 147\,MiB of Peer
    residency removes preload wait. Totals end at tail-KV reconstruction.
    The homogeneous, preemptible-copy model excludes detection,
    queuing, and interference.}
  \label{fig:design-pipeline}
  \Description{Four timelines compare token-block transfer, full-layer
    pipelining, hot preloads with demand fetch, and Peer-resident hot sets.}
\end{figure}

Figure~\ref{fig:design-pipeline} illustrates the depth tradeoff at one
fixed tail length: layer pipelining, hot preloads, and seven resident
hot sets reduce exposed wait from 172 to 83, 44, and 20\,ms.
This calibrated example does not validate the adaptive controller or
make seven layers appropriate for every checkpoint interval.
